\section{Kết luận \& Hướng phát triển}
\label{sec:ch8}

\subsection{Kết quả đạt được}

Dự án đã xây dựng hoàn chỉnh một DApp đặt chỗ tour du lịch chạy trên Ethereum Sepolia, đáp ứng đủ yêu cầu của Đề 36:

\begin{itemize}[leftmargin=1.5em]
    \item \textbf{Smart contract} \code{TourReservation} (Solidity \code{\^{}0.8.20}) hiện thực bốn hàm \code{createTour()}, \code{bookSeats()},
          \code{getAvailableSeats()}, \code{getMyBookings()}; lưu số vé theo ví bằng \code{mapping}; kiểm soát quyền điều hành bằng \code{onlyOperator};
          kiểm tra số chỗ dương, không vượt chỗ trống và thanh toán đúng phí bằng \code{require}; phát sự kiện \code{TourCreated}, \code{SeatsBooked}.
    \item \textbf{Frontend} React + Vite + web3.js tự nhận diện vai trò (\code{isOperator}) để tách giao diện Điều hành và Khách hàng,
          hiển thị giá vé, chỗ trống, vé đã đặt, đồng bộ dữ liệu sau mỗi giao dịch và khi đổi tài khoản.
    \item \textbf{Kết nối Web3} qua MetaMask (\code{eth\_requestAccounts}, \code{accountsChanged}, \code{chainChanged}); giao dịch được ký trong ví, phí gas do người gửi trả.
    \item \textbf{Kiểm thử thực nghiệm} trên Sepolia với ba ví (Điều hành, Khách hàng 1, Khách hàng 2) cho thấy các ràng buộc nghiệp vụ được đáp ứng
          và các tình huống lỗi được chặn ở cả ba lớp: giao diện, ví MetaMask và smart contract.
\end{itemize}

Qua dự án, có thể thấy lợi thế của kiến trúc DApp: quy tắc nghiệp vụ nằm trong mã hợp đồng nên minh bạch và không thể bị chỉnh sửa tùy tiện sau khi triển khai;
tiền của khách được hợp đồng giữ và mọi giao dịch đều tra cứu được công khai trên Etherscan.

\subsection{Hạn chế hiện tại}

\begin{itemize}[leftmargin=1.5em]
    \item \textbf{Chưa có cách rút tiền:} số ETH thu từ vé nằm trong hợp đồng nhưng không có hàm để Điều hành rút, nên tiền bị khóa trong hợp đồng.
    \item \textbf{Không hủy vé/hoàn tiền:} khách đã đặt không thể hủy.
    \item \textbf{Số vé theo ví chỉ xem được qua ví đó:} chưa có danh sách khách hàng hay lịch sử đặt chỗ trên giao diện.
    \item \textbf{Chưa có kiểm thử tự động} cho hợp đồng (hiện kiểm thử thủ công trên Sepolia).
\end{itemize}

\subsection{Hướng phát triển}

\begin{enumerate}[leftmargin=1.8em]
    \item \textbf{Hàm \code{withdraw()}:} cho phép riêng Điều hành rút ETH về ví của mình (\code{onlyOperator}), có thể chỉ cho phép sau ngày khởi hành tour.
          Nên dùng mẫu \emph{pull payment} và kiểm tra--cập nhật--chuyển tiền (\emph{checks-effects-interactions}) để tránh tấn công tái nhập (reentrancy).
    \item \textbf{Nhiều tour theo \code{tourId}:} định nghĩa \code{struct Tour \{ totalSeats; availableSeats; ticketPrice; departureTime; active; \}}, lưu bằng
          \code{mapping(uint256 => Tour)}, và số vé theo hai khóa \code{mapping(uint256 => mapping(address => uint256))}. Khi đó
          \code{getAvailableSeats(tourId)} và \code{getMyBookings(tourId)} nhận thêm tham số, giao diện có danh sách tour để chọn.
    \item \textbf{Hủy vé/hoàn tiền:} thêm hàm \code{cancelBooking(tourId, seats)} hoàn lại ETH theo chính sách (ví dụ hoàn đầy đủ trước hạn chót), trả lại chỗ trống
          và phát sự kiện tương ứng.
    \item \textbf{Lịch sử đặt chỗ từ sự kiện:} dùng \code{getPastEvents('SeatsBooked')} (đã có \code{customer} được \code{indexed}) để hiện danh sách đặt chỗ mà không cần lưu thêm dữ liệu trên chuỗi.
    \item \textbf{Kiểm thử tự động:} viết unit test bằng Hardhat hoặc Foundry cho mọi nhánh \code{require}, và chạy trên mạng cục bộ trước khi triển khai Sepolia.
    \item \textbf{Hoàn thiện triển khai:} cố định phiên bản web3.js trong \code{package.json}, triển khai giao diện lên dịch vụ lưu trữ tĩnh (Vercel/Netlify),
          và xác minh mã nguồn hợp đồng trên Etherscan để người dùng đọc được mã.
\end{enumerate}
